% Part: first-order-logic
% Chapter: models-theories
% Section: set-theory

\documentclass[../../../include/open-logic-section]{subfiles}

\begin{document}

\olfileid{fol}{mat}{set}

\olsection{سیٹاں دا نظریہ}

تقریباً ساری ریاضی سیٹاں دے نظریے وچ تشکیل دتی جا سکدی اے۔
ایس نظریے وچ ریاضی دی تشکیل لئی کئی کم کرن پینڈے نیں۔ سب توں
پہلاں تعلق~$\in$ لئی مسلّماں دا اک سیٹ چاہیدا اے۔ کجھ وکھرے
مسلّمی نظام بنائے گئے نیں، جنہاں وچ کدی کدی~$\in$ دیاں
خاصیتاں وی اک دوجے نال ٹکرا جاندیاں نیں۔ $\Log{ZFC}$ ناں دا
مسلّمی نظام، یعنی چناؤ دے مسلّمے نال زرمیلو--فرینکل سیٹ نظریہ،
انہاں وچ نمایاں اے: ایہ سب توں ودھ ورتیا تے پڑھیا گیا نظام اے،
کیونکہ ایہدے مسلّمے تقریباً اوہ ساریاں گلاں ثابت کرن لئی کافی
نکلدے نیں جیہڑیاں ریاضی دان ثابت کر سکن دی امید رکھدے نیں۔
پر ایہ قائم کرن توں پہلاں ایہ واضح کرن ضروری اے کہ ریاضی دان
جیہڑیاں گلاں کہنا چاہندے نیں، انہاں دا \emph{اظہار} کیویں
کیتا جا سکدا اے۔ شروع وچ ای زبان وچ کوئی !!{constant}s یا
!!{function}s نہیں، ایس لئی پہلی نظر وچ ایہ صاف نہیں کہ اسیں
کسے خاص سیٹ (مثلاً $\emptyset$ یا $\Nat$) بارے گل کیویں کریے،
سیٹاں اُتے عملاں (مثلاً $X \cup Y$ تے $\Pow{X}$) بارے کیویں
گل کریے، یا سیٹاں توں علاوہ ہور چیزاں، جیویں تعلقاں تے
فنکشناں، نوں شامل کرن والیاں بناوٹاں بارے گل کیویں کریے۔

شروع وچ ”دا رکن ہونا“ واحد تعلق نہیں جیہدے نال ساڈی دلچسپی اے:
”دا ذیلی سیٹ ہونا“ وی تقریباً اونا ای اہم لگدا اے۔ پر اسیں
”دا ذیلی سیٹ ہونا“ نوں ”دا رکن ہونا“ دے ذریعے \emph{تعریف}
کر سکدے آں۔ ایہدے لئی سیٹ نظریے دی زبان وچ اک
!!a{formula}~$!A(x, y)$ لبھنا اے جیہڑا سیٹاں دی جوڑی
$\tuple{X, Y}$ لئی جے تے صرف جے پورا اترے جدوں
$X \subseteq Y$ ہووے۔ پر $X$، $Y$ دا ذیلی سیٹ اوہدوں تے
صرف اوہدوں اے جدوں~$X$ دے سارے !!{element}s،
~$Y$ دے وی !!{element}s ہون۔ ایس لئی اسیں
$\subseteq$ نوں ہیٹھلے فارمولے نال تعریف کر سکدے آں:
\[
\lforall[z][(z \in x \lif z \in y)]
\]
ہن جتھے وی کسے فارمولے وچ تعلق~$\subseteq$ ورتنا ہووے،
اسیں اوہدی تھاں ایہ فارمولا ورت سکدے آں ($x$ تے $y$ نوں
مناسب طرح بدل کے، تے لوڑ پئے تاں پابند متغیر~$z$ دا ناں
بدل کے)۔ مثال لئی، سیٹاں دے عنصراں نال تعیّن دا مطلب ایہ
اے کہ جے کوئی دو سیٹ~$x$ تے $y$ اک دوجے وچ شامل ہون،
تاں $x$ تے $y$ اکّو سیٹ ہون گے۔ ایہ گل
$\lforall[x][\lforall[y][((x \subseteq y \land y
    \subseteq x) \lif x = y)]]$ نال ظاہر کیتی جا سکدی اے،
یا جے $\subseteq$ دی تھاں اُپرلی تعریف رکھیے، تاں
\[
\lforall[x][\lforall[y][((\lforall[z][(z \in x \lif z \in y)] \land
    \lforall[z][(z \in y \lif z \in x)]) \lif x = y)]].
\]
ایہ حقیقت وچ $\Log{ZFC}$ دے مسلّماں وچوں اک اے، یعنی
”عنصراں نال تعیّن دا مسلّمہ“۔

$\emptyset$ لئی کوئی !!{constant} نہیں، پر ”$x$ خالی اے“
نوں $\lnot \lexists[y][y \in x]$ نال ظاہر کیتا جا سکدا اے۔
فیر ”$\emptyset$ موجود اے“ !!{sentence}~$\lexists[x][\lnot \lexists[y][y \in
    x]]$ بن جاندا اے۔ ایہ~$\Log{ZFC}$ دا اک ہور مسلّمہ اے۔
(دھیان دیو کہ عنصراں نال تعیّن دا مسلّمہ ایہ لازم کردا اے
کہ خالی سیٹ صرف اک ای اے۔) سیٹ نظریے دی زبان وچ جدوں وی
$\emptyset$ بارے گل کرنی ہووے، اسیں ایویں لکھاں گے:
”اک سیٹ اے جیہڑا خالی اے تے \dots“۔ مثال دے طور اُتے،
ایہ ظاہر کرن لئی کہ $\emptyset$ ہر سیٹ دا ذیلی سیٹ اے،
اسیں لکھ سکدے آں
\[
\lexists[x][(\lnot\lexists[y][y \in x] \land \lforall[z][x \subseteq
    z])]
\]
جتھے بے شک $x \subseteq z$ نوں وی اوہدی تعریف نال
بدلنا پئے گا۔

سیٹاں اُتے عملاں، جیویں $X \cup Y$ تے $\Pow{X}$، بارے گل
کرن لئی وی اسیں ایہو جیہا طریقہ ورتنا اے۔ سیٹ نظریے دی
زبان وچ فنکشن دیاں علامتاں نہیں، پر فنکشنی تعلق
$X \cup Y = Z$ تے $\Pow{X} = Y$ نوں اسیں ایویں ظاہر
کر سکدے آں:
\begin{align*}
& \lforall[u][((u \in x \lor u \in y) \liff u \in z)]\\
& \lforall[u][(u \subseteq x \liff u \in y )]
\end{align*}
کیونکہ $X \cup Y$ دے !!{element}s عین اوہ سیٹ نیں جیہڑے
یا~$X$ دے !!{element}s ہون یا~$Y$ دے !!{element}s،
تے $\Pow{X}$ دے !!{element}s عین~$X$ دے ذیلی سیٹ نیں۔
پر ایہدے نال اسیں $x \cup y$ یا $\Pow{x}$ نوں ٹرماں وانگ
نہیں ورت سکدے؛ اسیں صرف اوہ پورے !!{formula}s ورت سکدے
آں جیہڑے $X \cup Y = Z$ تے $\Pow{X} = Y$ تعلقاں دی
تعریف کردے نیں۔ حقیقت وچ اسیں اجے ایہ وی نہیں جاندے کہ
ایہ تعلق کدے پورے وی اتردے نیں یا نہیں؛ یعنی ایہ نہیں
جاندے کہ اتحاد سیٹ تے قوت سیٹ ہمیشہ موجود ہوندے نیں۔
مثال لئی !!{sentence}
$\lforall[x][\lexists[y][\Pow{x} = y]]$،
~$\Log{ZFC}$ دا اک ہور مسلّمہ اے (قوت سیٹ دا مسلّمہ)۔

ہن ترتیب وار جوڑیاں یا فنکشناں بارے کیویں گل کریے؟ ایتھے
سانوں دسنا پئے گا کہ ترتیب وار جوڑیاں تے فنکشناں نوں
خاص قسم دے سیٹ کیویں سمجھیا جا سکدا اے۔ ترتیب وار جوڑی
$\tuple{x, y}$ دی اک تعریف سیٹ $\{\{x\}, \{x, y\}\}$
اے۔ پر پہلاں وانگ اسیں ایس سیٹ دا ناں رکھن والا
!!a{function} شامل نہیں کر سکدے؛ اسیں صرف تعلق
$\tuple{x, y} = z$، یعنی $\{\{x\}, \{x, y\}\} = z$،
دی تعریف کر سکدے آں:
\[
\lforall[u][(u \in z \liff (\lforall[v][(v \in u \liff v = x)] \lor
  \lforall[v][(v \in u \liff (v = x \lor v = y))]))]
\]
ایہ کہندا اے کہ !!{element}s~$u$، جیہڑے~$z$ وچ نیں، عین اوہ سیٹ
نیں جنہاں دا اکلوتا !!{element} $x$ اے، یا جنہاں دے
صرف !!{element}s $x$ تے~$y$ نیں؛ دوجے لفظاں وچ،
اوہ سیٹ جیہڑے یا $\{x\}$ دے برابر نیں یا
$\{x, y\}$ دے برابر نیں۔ ایہ بناوٹ ملن توں بعد اسیں
ہور گل وی کہہ سکدے آں، مثلاً $X \times Y = Z$:
\[
\lforall[z][(z \in Z \liff \lexists[x][\lexists[y][(x \in
        X \land y \in Y \land \tuple{x, y} = z)]])]
\]

فنکشن $f \colon X \to Y$ نوں تعلق $f(x)
= y$ دے روپ وچ، یعنی جوڑیاں دے سیٹ
$\Setabs{\tuple{x,y}}{f(x) = y}$ دے روپ وچ سمجھیا جا سکدا اے۔
فیر اسیں کہہ سکدے آں کہ سیٹ~$f$، $X$ توں $Y$ ول فنکشن
اے جے (الف) ایہ $\subseteq X \times Y$ تعلق اے، (ب) ایہ
کُل اے، یعنی ہر $x \in X$ لئی کوئی $y \in Y$ اے
جیہدے لئی $\tuple{x, y} \in f$ ہووے، تے (ج) ایہ
فنکشنی اے، یعنی جدوں وی $\tuple{x, y}, \tuple{x, y'} \in f$
ہوون، $y = y'$ ہووے (کیونکہ اکّو انپٹ لئی فنکشن دی
قدر اکّو ہونی چاہیدی اے)۔ ایس لئی ”$f$، $X$ توں $Y$ ول
فنکشن اے“ ایویں لکھیا جا سکدا اے:
% PNB-SYN-008: فنکشن دے بیان وچ باہرلے مشروط بیان دی رہ گئی قوس بند کیتی اے۔
\begin{align*}
 \lforall[u][(u \in f \lif {}] & \lexists[x][\lexists[y][(x \in X \land y \in
      Y \land \tuple{x, y} = u)]]) \land {}\\
 \lforall[x][(x \in X \lif {}] &
      (\lexists[y][(y \in Y \land \mathrm{maps}(f, x, y))] \land {}\\
& (\lforall[y][\lforall[y'][((\mathrm{maps}(f, x, y) \land
    \mathrm{maps}(f, x, y')) \lif y = y')]])))
\end{align*}
جتھے $\mathrm{maps}(f, x, y)$، $\lexists[v][(v \in f \land
  \tuple{x, y} = v)]$ دا مخفف اے (ایہ !!{formula}،
”$f(x) = y$“ دا اظہار کردا اے)۔

ہن ایہ ظاہر کرن وی اوکھا نہیں کہ $f\colon X \to Y$
!!{injective} اے، مثلاً:
\begin{multline*}
 f \colon X \to Y \land \lforall[x][\lforall[x'][((x \in X \land x' \in
     X \land {}]] \\
 \lexists[y][(\mathrm{maps}(f, x, y) \land \mathrm{maps}(f,
        x', y))]) \lif x = x')
\end{multline*}
فنکشن~$f\colon X \to Y$ !!{injective} اے جے تے صرف
جے، جدوں وی $f$، $x, x' \in X$ نوں اکّو~$y$
ول بھیجے، $x = x'$ ہووے۔ جے اسیں ایس فارمولے
دا مخفف $\mathrm{inj}(f, X, Y)$ رکھ لئیے، تاں ہن
سیٹ نظریے دی زبان وچ کانتور دے قضیے ورگی اہم گل
وی کہہ سکدے آں: $\Pow{X}$ توں $X$ ول کوئی
!!{injective} فنکشن نہیں:
\[
\lforall[X][\lforall[Y][(\Pow{X} = Y \lif
    \lnot\lexists[f][\mathrm{inj}(f, Y, X)])]]
\]

کسے نوں لگ سکدا اے کہ سیٹ نظریے نوں اک ہور مسلّمہ چاہیدا
اے جیہڑا ہر وصف لئی اک سیٹ دی موجودگی یقینی بناوے۔
% PNB-SEM-025: ایتھے دکھائے گئے جملے لئی ہور آزاد متغیر نہیں؛ پیرامیٹر والے عمومی خاکے دی بندش ہیٹھاں واضح اے۔
جے $!A(x)$ سیٹ نظریے دا اوہ فارمولا اے جیہدے وچ
صرف متغیر~$x$ آزاد ہو سکدا اے، تاں اسیں !!{sentence}
\[
\lexists[y][\lforall[x][(x \in y \liff !A(x))]].
\]
اُتے غور کر سکدے آں۔ ایہ !!{sentence} کہندا اے کہ
سیٹ~$y$ موجود اے جیہدے !!{element}s اوہ سارے تے
صرف اوہ $x$ نیں جیہڑے~$!A(x)$ نوں پورا کردے نیں۔
ایس خاکے نوں ”خاصیت توں سیٹ بنان دا اصول“ کہندے نیں۔
ایہ بڑا کارآمد لگدا اے؛ پر بدقسمتی نال ایہ غیر موافق اے۔
$!A(x) \ident \lnot x \in x$ رکھو؛ فیر خاصیت توں
سیٹ بنان دا اصول کہندا اے:
\[
\lexists[y][\lforall[x][(x \in y \liff x \notin x)]],
\]
یعنی اوہ سارے سیٹاں دا اک سیٹ موجود اے جیہڑے
اپنے آپ دے !!{element}s نہیں۔ ایسا کوئی سیٹ موجود
نہیں ہو سکدا---ایہ رسل دا پیراڈوکس اے۔ حقیقت وچ
$\Log{ZFC}$ ایس اصول دا اک محدود---تے موافق---روپ
رکھدا اے، جیہنوں وکھ کرن دا اصول کہندے نیں:
% PNB-SYN-008: وکھ کرن دے خاکے وچ باہرلی دو طرفہ شرط دی رہ گئی قوس بند کیتی اے۔
\[
\lforall[z][\lexists[y][\lforall[x][(x \in y \liff (x \in z \land
      !A(x)))]]].
\]
عمومی خاکے وچ وادھو آزاد متغیراں نوں پیرامیٹر من کے پورے
بیان نوں انہاں اُتے ہمہ گیر مسوّراں نال بند کیتا جاندا اے؛
باہرلے مسوّراں دے متغیر نام تازہ رکھے جاندے نیں۔ ”موافق“
ایتھے وکھ کرن دے اصول بارے اے، پورے سیٹ نظریے دی بےتضادی
دا ثبوت نہیں۔

\begin{prob}
اک ایداں دا !!a{derivation} دے کے دکھاؤ جیہدے نال ثابت ہووے
کہ خاصیت توں سیٹ بنان دا اصول غیر موافق اے، یعنی
\[
\lexists[y][\lforall[x][(x \in y \liff x \notin x)]] \Proves \lfalse.
\]
پہلاں ایہ ثابت کرنا مددگار ہو سکدا اے کہ
$(A \lif \lnot A) \land (\lnot A \lif A)
\Proves \lfalse$۔
\end{prob}
\end{document}
